\section{Introduction}
\label{sec:introduction}

The Boltzmann--Grad limit describes a gas whose interaction range tends
to zero while the collision frequency remains of order one. In three
dimensions the activity is $\mu_\varepsilon=\varepsilon^{-2}$.
The microscopic dynamics is Newtonian, whereas the limiting collision
operator is determined by the actual two-body scattering of the pair
potential. We prove this limit for stable finite-range radial potentials
on every finite interval where the Boltzmann solution has the regularity
and Gaussian bound specified below.

Write $z=(x,v)\in\R^3\times\R^3$.
For a phase-space function $h$ and $\gamma>0$, write
\begin{equation}\label{intro:norm}
 \|h\|_{\mathrm{Bol},\gamma}
 :=\sum_{k\in\Z^3}\;
 \sup_{\substack{|x-k|\le1\\v\in\R^3}}
 e^{\gamma|v|^2}|h(x,v)|.
\end{equation}
For a vector-valued function we use its Euclidean length inside this
definition. The precise potential class, grand-canonical law, Newton
equations, and scattering operator are given in
Assumption~\ref{pre:potential} and Equations~\eqref{pre:model},
\eqref{pre:dynamics}, and \eqref{pre:collision}.
In brief, the potential is radial,
vanishes beyond unit distance, and is $C^2$, with either a finite value
or a repulsive singularity at the origin. Its only sign condition is
thermodynamic stability, in the classical sense of~\cite{Ruelle69}.
The initial law is the product activity $f_0$
conditioned by pairwise exclusion at distance $\varepsilon$, with its
exact grand-canonical normalization.
For each fixed integer $s\ge1$, the scaled factorial density
$F_s^\varepsilon(t)$ is characterized by
\begin{equation}\label{intro:factorial-test}
 \int a(Z_s)F_s^\varepsilon(t,Z_s)\,dZ_s
 =\mu_\varepsilon^{-s}\EE
   \sum_{\substack{1\le i_1,\ldots,i_s\le N\\
                    i_1,\ldots,i_s\ \mathrm{distinct}}}
   a(z_{i_1}(t),\ldots,z_{i_s}(t))
\end{equation}
for bounded measurable tests $a$ on $(\R^3\times\R^3)^s$.
The sum runs over ordered distinct labels and is zero when $N<s$;
$Z_s=(z_1,\ldots,z_s)$ denotes an ordered phase-space configuration.

\begin{maintheorem}[Regular-lifespan kinetic limit]\label{thm:main}
Let $\Phi$ satisfy Assumption~\ref{pre:potential}, and let $\beta>0$.
Let $f_0\ge0$ be a $C^1$ probability density on $\R^3\times\R^3$ such that
\begin{equation}\label{intro:data-bound}
 \|f_0\|_{\mathrm{Bol},2\beta}
 +\|\nabla_xf_0\|_{\mathrm{Bol},2\beta}<\infty.
\end{equation}
Fix $0<T<\infty$ for which the equation
\begin{equation}\label{intro:boltzmann}
 (\partial_t+v\cdot\nabla_x)f=Q_\Phi(f,f),
 \qquad f|_{t=0}=f_0,
\end{equation}
has a nonnegative classical probability-density solution satisfying
\begin{equation}\label{intro:solution-bound}
 \sup_{0\le t\le T}\sup_{x,v}
 e^{2\beta|v|^2}f(t,x,v)<\infty.
\end{equation}
Start the exact Newtonian process from the grand-canonical law
\eqref{pre:model}. For its scaled factorial densities
$F_s^\varepsilon$, defined in \eqref{pre:factorial}, every fixed
integer $s\ge1$ satisfies
\begin{equation}\label{intro:chaos}
 \lim_{\varepsilon\downarrow0}\;
 \sup_{0\le t\le T}
 \bigl\|F_s^\varepsilon(t)-f(t)^{\otimes s}\bigr\|_
 {L^1((\R^3\times\R^3)^s)}=0.
\end{equation}
\end{maintheorem}

The potential and $T$ are fixed independently of $\varepsilon$.
The convergence modulus may depend on the potential and on the bounds
in \eqref{intro:data-bound}--\eqref{intro:solution-bound}.
Attractive wells, long scattering delays, and noninjective scattering
maps are permitted. The exclusion belongs solely to the initial law.
The Newtonian evolution retains simultaneous interactions,
recollisions, and any clusters it creates.

\begin{maincorollary}[Empirical observables]\label{cor:empirical}
Under the hypotheses of Theorem~\ref{thm:main}, let
$a\in C_c^\infty(\R^3\times\R^3;\R)$ and $\eta>0$. Then
\begin{equation}\label{intro:empirical}
 \lim_{\varepsilon\downarrow0}\;
 \sup_{0\le t\le T}
 \PP\!\left(
 \left|\mu_\varepsilon^{-1}\sum_{i=1}^{N}a(z_i(t))
       -\int a(z)f(t,z)\,dz\right|>\eta\right)=0.
\end{equation}
\end{maincorollary}

The supremum in \eqref{intro:empirical} is over deterministic
observation times. The conclusion follows from the first two factorial
marginals and conservation of particle number; see
Section~\ref{sec:closure}.

\subsection{History and scope}

The Boltzmann--Grad scaling goes back to Grad~\cite{Grad49,Grad58}.
Lanford established the hard-sphere limit on a short time
interval~\cite{Lanford75}, and King's thesis treated positive
finite-range potentials~\cite{King75}.
Gallagher, Saint-Raymond, and Texier gave a detailed geometric analysis
of recollisions for hard spheres and a class of repulsive potentials,
including a convergence rate in the hard-sphere case~\cite{GST14}.
Pulvirenti, Saffirio, and Simonella proved a short-time theorem for the
stable radial class used here; see their Section~8, Hypothesis~1$'$ and
Theorem~1$'$~\cite{PSS14}. Their treatment permits attractive parts and
does not require a globally single-valued differential cross-section.

The time restriction in a collision-tree expansion reflects the growth
of correlated histories. Global-time results were obtained for
sufficiently rare, spatially dispersive gases in vacuum by Illner and
Pulvirenti~\cite{IllnerPulvirenti86,Pulvirenti87,IllnerPulvirenti89};
the last reference corrects and improves the earlier results.
For general data, extending a short-time limit requires control of the
correlations already present at the next starting time.
Pulvirenti and Simonella quantified the correlation errors remaining
after independent factors are removed~\cite{PulvirentiSimonella2017}.
Bodineau, Gallagher, Saint-Raymond, and Simonella developed a cluster
expansion directly on physical trajectories~\cite{BGSSCluster2022},
distinguishing genuine collisions from geometric overlaps between
independently evolved components.

Deng, Hani, and Ma proved the hard-sphere limit throughout every regular
finite kinetic interval~\cite[Theorem~1]{DHM25}. Their layered expansion
retains the earlier histories of correlated remainders while extracting
the factors approximated by the kinetic solution. This strategy has an
antecedent in the long-time wave-kinetic work of Deng and
Hani~\cite{DengHani2024}. The exposition~\cite{BGSS26} explains both the
accumulated correlations and the loss of topology that prevent
restarting a short-time chaos theorem as a black box.

Smooth potentials require additional control because encounters last
for positive time and several particles can interact simultaneously.
Deng, Hani, and Ma discuss this extension as a further problem
in~\cite[Section~1.4.2]{DHM25}.
The present whole-component expansion retains the potential energy of
these interactions. Its stability and causal exposure estimates are
proved below; they are not supplied by the hard-sphere convergence
theorem. Similarly, the assumed classical kinetic solution is an input
to the theorem, while its additional spatial envelopes are derived in
Section~\ref{sec:preliminaries} by adapting the argument of
\cite[Proposition~A.1]{DHM25}.

Theorem~\ref{thm:main} proves the regular-lifespan Boltzmann--Grad
statement for the full stable radial class. Its lifespan assumption
concerns the kinetic solution itself and imposes no small-data or
near-equilibrium condition. The conclusion is convergence of all fixed
factorial marginals in full phase-space $L^1$.

\subsection{Structure of the proof}

The argument uses a finite number of macroscopic layers of length
$b=T/L$, each divided into a much finer mesh. The fine cells permit an
exact finite expansion by whole interaction components. The coarse
layers permit centering against the known Boltzmann solution. Several
features of the construction are useful beyond ordinary binary
collision trees.
Throughout the proof, $P=1+|\log\varepsilon|$; polynomial cutoffs mean
fixed powers of $P$, chosen after the finite number $L$ of layers.

\paragraph{Exact component identities.}
The expansion groups particles into their entire instantaneous
interaction components at every cut. A finite inclusion-exclusion
over successive generations of components represents each coefficient
by independent subsystem flows and contact indicators. This preserves
the potential energy needed for estimates even when a component has
many internal encounters. Centering acts only on isolated singleton
slots. Its cancellation requires separation of every displayed
trajectory, including the detached trajectories in the signed terms.

\paragraph{One energy budget across all layers.}
Each formal particle lifetime contributes an independent input once
and departs once. Stability applied to intact departed components gives
a single Gaussian budget for the whole history. Repeatedly carried
velocities can then cost $L^{p/2}$ for a history with $p$ labels, while
the time factors contribute $(T/L)^{p-h_0}$ for $h_0$ top roots.
The $h_0$ roots are the labels at the upper end from which the
backward history starts. After bounding their contribution, the resulting
base
$CT/\sqrt L$ is small for a sufficiently large fixed $L$. This is the
mechanism that passes beyond the convergence interval of one
collision-tree series.

\paragraph{Causal exposure of defects.}
The first estimates integrate all positions and velocities at the
ends of a fine time cell. They bound histories with crowded groups or
several additions in one cell. The remaining histories can be represented
by trees whose edges follow the order in which particles are added.
If extra contacts affect a positive fraction of the particle lives,
a counting argument selects many disjoint trees of bounded size.
Their fresh scattering parameters give small angular sets for contacts
with paths that are independent of those parameters.
To multiply these small factors, both outputs of an unexposed genuine
scattering must remain hidden, whereas a virtual crossing hides only
the new branch; Figure~\ref{fig:exposure} explains the distinction.
The trees used to choose integration coordinates need not coincide
with the families used in the signed cancellation identity.

\begin{figure}[tbp]
\centering
\begin{tikzpicture}[x=0.94cm,y=0.94cm,
 visible/.style={draw=black!75,line width=1.05pt},
 hidden/.style={draw=blue!65!black,line width=1.2pt},
 every node/.style={font=\small},
 exposure point/.style={circle,fill=red!65!black,inner sep=2pt}]
\draw[-{Latex[length=2mm]},black!55] (-2.6,3.1)--(-2.6,-.05);
\node[rotate=90,anchor=south,font=\footnotesize] at (-2.72,1.5)
 {backward time};
\begin{scope}
 \node[font=\small\bfseries] at (.15,3.65) {Genuine scattering};
 \fill[blue!4] (-1.65,-.15) rectangle (1.6,1.82);
 \draw[blue!35,dashed,rounded corners=3pt]
  (-1.65,-.15) rectangle (1.6,1.82);
 \draw[visible] (0,3.15)--(0,2.1);
 \draw[hidden] (0,2.1)--(-1.18,0);
 \draw[hidden] (0,2.1)--(1.18,0);
 \fill[black] (0,2.1) circle (1.7pt);
 \node[anchor=east,font=\footnotesize] at (-.13,2.32) {launch};
 \draw[visible] (2.2,3.15) .. controls (2.2,1.85) and (1.75,1.05)
  .. (.787,.7);
 \node[exposure point] at (.787,.7) {};
 \node[anchor=west,font=\footnotesize] at (1.5,.38) {first exposure};
 \node[align=center,font=\footnotesize] at (-1.18,-.48) {new label};
 \node[align=center,font=\footnotesize] at (1.18,-.48) {continued parent};
 \node[blue!65!black,font=\footnotesize] at (.05,-.98)
  {both descendant branches hidden};
\end{scope}
\begin{scope}[xshift=7cm]
 \node[font=\small\bfseries] at (.25,3.65) {Passage or virtual crossing};
 \fill[blue!4] (.18,-.15) rectangle (1.65,1.82);
 \draw[blue!35,dashed,rounded corners=3pt]
  (.18,-.15) rectangle (1.65,1.82);
 \draw[visible] (0,3.15)--(0,0);
 \draw[hidden] (0,2.1)--(1.18,0);
 \fill[black] (0,2.1) circle (1.7pt);
 \node[anchor=east,font=\footnotesize] at (-.13,2.32) {launch};
 \draw[visible] (2.2,3.15) .. controls (2.2,1.85) and (1.75,1.05)
  .. (.787,.7);
 \node[exposure point] at (.787,.7) {};
 \node[anchor=east,font=\footnotesize] at (-.06,-.48) {visible parent};
 \node[anchor=west,font=\footnotesize] at (.9,-.48) {new label};
 \node[blue!65!black,font=\footnotesize] at (.3,-.98)
  {only the new branch hidden};
\end{scope}
\end{tikzpicture}
\caption{The auxiliary exposure used to multiply contact savings.
Blue branches carry an unrevealed launch parameter. A genuine
scattering changes the continued parent's velocity, so both descendant
branches must be hidden. Passage leaves the parent unchanged.
Until first exposure, the black trajectories can be computed without
the hidden parameter. The construction is applied to exact signed
histories; no particles are deleted from the microscopic dynamics.
The drawing shows one selected tree. The full construction also omits
added groups whose positions have not yet been fixed relative to the
visible particles, including groups not selected for exposure;
Proposition~\ref{exp:exposure} gives the complete construction.}
\label{fig:exposure}
\end{figure}


Section~\ref{sec:preliminaries} establishes the flow, initial estimates,
scattering facts, and uniform smooth-input bounds.
Section~\ref{sec:cell-expansion} constructs the exact finite expansion.
Section~\ref{sec:centering} gives the centered identities and the
short-layer approximation to the kinetic solution.
Sections~\ref{sec:history-estimates} and~\ref{sec:exposure} prove the
finite-history measure estimates and the exposure bound.
Section~\ref{sec:closure} chooses the thresholds, removes every cutoff,
and proves Theorem~\ref{thm:main} and Corollary~\ref{cor:empirical}.
